\subsection*{6.5. Picard--Lefschetz 公式的回顾 (XV 3.4)}

设 $S$ 为严格 Hensel 赋值环的谱, 泛点为 $\eta$, 几何泛点为 $\bar\eta$, 闭点为 $s$. 设
\[
\rho:X\longrightarrow S
\]
为固有态射, 其中 $X$ 为 $n$ 维连通正则概形. 假设 $\rho$ 除特殊纤维中的唯一一点 $x$ 外处处光滑, 且 $x$ 为该纤维的普通二次奇点.

以 $I$ 表示 $S$ 在 $s$ 处的惯性群. 由 XV 3.4, 当 $i\ne n-1$ 时, $I$ 在 $H^i(X_{\bar\eta},\mathbb Q_\ell)$ 上的作用平凡, 而它在 $H^{n-1}(X_{\bar\eta},\mathbb Q_\ell)$ 上的作用描述如下.

\subsubsection*{6.5.1.}
若 $n=2m+1$, 则 $I$ 的作用通过一个同构于 $\mathbb Z/2\mathbb Z$ 的商群分解 (若 $\operatorname{char}(k)\ne2$, 则仅有一个这样的商群). 存在元素
\[
\delta\in H^{n-1}(X_{\bar\eta},\mathbb Q_\ell(m)),
\]
相差符号时唯一确定, 称为\underline{$s$ 处的消灭圈}, 使得
\begin{equation}\tag{6.5.1.1}
\delta\cdot\delta=(-1)^m2
\quad\text{在}\quad
H^{2n-2}(X_{\bar\eta},\mathbb Q_\ell(n-1))\simeq\mathbb Q_\ell.
\end{equation}
$\mathbb Z/2\mathbb Z$ 的非平凡元素 $\sigma$ 在 $H^{n-1}(X_{\bar\eta},\mathbb Q_\ell)$ 上按反射
\begin{equation}\tag{6.5.1.2}
\sigma x=x-(-1)^m(x\cdot\delta)\delta.
\end{equation}
\footnote{印本将反射公式编号为 (6.5.2.2), 但它是 6.5.1 的第二个公式, 且后接 (6.5.2.1).}
作用. 这里 $x\cdot\delta$ 属于 $H^{2n-2}(X_{\bar\eta},\mathbb Q_\ell(m))\simeq\mathbb Q_\ell(-m)$, 所以 $(x\cdot\delta)\delta$ 定义 $H^{n-1}(X_{\bar\eta},\mathbb Q_\ell)$ 的一个元素.

\subsubsection*{6.5.2.}
若 $n=2m+2$, 则 $I$ 的作用通过其典范商群 $\mathbb Z_\ell(1)$ 分解. 存在元素
\[
\delta\in H^{n-1}(X_{\bar\eta},\mathbb Q_\ell(m)),
\]
相差符号时典范确定, 称为\underline{$s$ 处的消灭圈}, 使得 $\sigma\in\mathbb Z_\ell(1)$ 按错切变换作用于 $x\in H^{n-1}(X_{\bar\eta},\mathbb Q_\ell)$:
\begin{equation}\tag{6.5.2.1}
\sigma(x)=x-(-1)^m\sigma\,(x\cdot\delta)\delta.
\end{equation}
事实上, $x\cdot\delta$ 属于 $H^{2n-2}(X_{\bar\eta},\mathbb Q_\ell(m))\simeq\mathbb Q_\ell(-m-1)$. 因此与 $\sigma\in\mathbb Z_\ell(1)$ 相乘得到元素 $\sigma(x\cdot\delta)\in\mathbb Q_\ell(-m)$, 它与 $\delta\in H^{n-1}(X_{\bar\eta},\mathbb Q_\ell(m))$ 的乘积属于 $H^{n-1}(X_{\bar\eta},\mathbb Q_\ell)$.

回到整体情形.

\subsubsection*{定理 6.6.}
设 $X$ 为代数闭域 $k$ 上的 $n$ 维光滑连通射影簇. 假设 $n$ 为偶数, 或 $\operatorname{char}(k)\ne2$. 设 $D$ 为 $X$ 的一个 Lefschetz 线束, $\eta$ 为其泛点, $\bar\eta$ 为 $\eta$ 上的一个几何点, $s_1,\ldots,s_N$ 为 $D\cap\check X$ 的各点.

对每个 $i$, 选取 $\bar\eta$ 与 $S_i$ 的几何泛点之间的同构 $\ell_i$, 其中后者是 $D$ 在 $s_i$ 处的严格 Hensel 化. 经由 $\ell_i$, 以 $I_i$ 表示 $s_i$ 处作用于 $H^{n-1}(X_{\bar\eta},\mathbb Q_\ell(j))$ 的惯性群, 以 $\delta_i$ 表示 6.5 定义的 $s_i$ 处的消灭圈.\footnote{印本在定义 $\delta_i$ 之前, 称 $I_i$ 为“$\delta_i$ 的”惯性群. 标准构造是 $s_i$ 处的惯性群, 通过所选道路 $\ell_i$ 嵌入.}
因而
\[
\delta_i\in H^{n-1}(X_{\bar\eta},\mathbb Q_\ell(m)),
\qquad
m=\begin{cases}
(n-1)/2,& n\text{ 为奇数},\\
n/2-1,& n\text{ 为偶数}.
\end{cases}
\]
\footnote{印本在此公式中略去了一般纤维的下标; 前面的作用以及 6.5 的定义都将每个 $\delta_i$ 置于 $X_{\bar\eta}$ 的上同调中.}
选取 $\ell_i$, 使子群
\[
I_i\subset
\pi=\pi_1^{\mathrm{mod}}(D-D\cap\check X,\bar\eta)
\]
在拓扑意义下生成 $\pi$, 由 SGA 1, XIII 2.10 可知这是可行的. 则:

\medskip
\noindent
\textup{6.6.1.} 若 $X$ 满足条件 (LV), 则消灭上同调空间
\[
E^{n-1}(X_{\bar\eta},\mathbb Q_\ell(m))
\stackrel{\mathrm{def}}{=}
H^{n-1}(X,\mathbb Q_\ell(m))^\perp
\subset H^{n-1}(X_{\bar\eta},\mathbb Q_\ell(m))
\]
(参见 5.4.4) 由消灭圈 $\delta_i$, $i=1,\ldots,N$ 生成.

\medskip
\noindent
\textup{6.6.2.} 消灭圈 $\delta_i$ 在整体单值群 $\pi_1(D-D\cap\check X)$ 的作用下彼此共轭, 允许相差符号.

\subsubsection*{证明.}
对于 6.6.1, Picard--Lefschetz 公式 6.5 和 $I_i$ 生成 $\pi=\pi_1^{\mathrm{mod}}(D-D\cap\check X)$ 这一事实给出
\begin{equation}\tag{6.6.3}
H^{n-1}(X_{\bar\eta},\mathbb Q_\ell(m))^\pi
=\bigl\langle\delta_1,\ldots,\delta_N\bigr\rangle^\perp
\quad\text{在 }H^{n-1}(X_{\bar\eta},\mathbb Q_\ell(m)).
\end{equation}
由 5.6.10,
\begin{equation}\tag{6.6.4}
H^{n-1}(X_{\bar\eta},\mathbb Q_\ell(m))^\pi
=\operatorname{Im}\bigl(
H^{n-1}(X,\mathbb Q_\ell(m))
\longrightarrow H^{n-1}(X_{\bar\eta},\mathbb Q_\ell(m))
\bigr).
\end{equation}
在 (6.6.3) 中取正交补, 即得 6.6.1.

对于 6.6.2, 只须注意, 消灭圈 $\delta_i$ 相差符号时由 $s_i$ 处惯性群的表示确定. 应用 6.1.2, 它给出复合惯性同态在 $\pi_1^{\mathrm{mod}}(\check{\mathbb P}^{r}-\check X,\bar\eta)$ 中的共轭性:
\begin{center}
\begin{tikzcd}[column sep=huge,row sep=huge]
\displaystyle\prod_{q\ne\operatorname{char}(k)}\mathbb Z_q(1) \arrow[r]
  & \pi_1^{\mathrm{mod}}(D-D\cap\check X,\bar\eta) \arrow[d] & {}\\
{}
  & \pi_1^{\mathrm{mod}}(\check{\mathbb P}^{r}-\check X,\bar\eta) & .
\end{tikzcd}
\end{center}
\begin{editorialnote}
原印本及英译底本在 6.6.1 的证明中将 $\pi$ 写成通常基本群; 此处沿用定理中已定义的驯基本群, 因为 6.1.10 说明表示通过该商分解, 并不说明两个基本群相等. 在 6.6.2 的惯性图中, 底本仅写 $\mathbb Z_\ell(1)$; 这里应用 6.1.2 时应取完整驯惯性群 $\prod_{q\ne\operatorname{char}(k)}\mathbb Z_q(1)$. 例如当 $n$ 为奇数而 $\ell\ne2$ 时, 反射的二阶商不能由其 $\ell$-部分检测. 此处据 6.1.2 的陈述订正.
\end{editorialnote}

因而, 由这些同态确定的消灭圈, 在 $\pi_1^{\mathrm{mod}}(\check{\mathbb P}^{r}-\check X,\bar\eta)$ 的元素作用下彼此共轭, 允许相差符号. 最后, 6.1.1 给出满射
\[
\pi_1^{\mathrm{mod}}(D-D\cap\check X)
\longrightarrow
\pi_1^{\mathrm{mod}}(\check{\mathbb P}^{r}-\check X).
\]

